\chapter{Getting started}\label{chapter-2}
\Needspace{10\baselineskip}
\section*{Obtain the source}
\pdfbookmark[1]{Obtain the source}{chapter-2-s1}
\par\noindent\begin{minipage}{\linewidth}
The project repository is \href{https://github.com/neerajjayesh/GridSec-project}{neerajjayesh/GridSec-project}.
\begin{lstlisting}[title={TEXT}]
git clone https://github.com/neerajjayesh/GridSec-project.git
cd GridSec-project
\end{lstlisting}
\end{minipage}\par
If you already have the source, open the directory containing \code{\seqsplit{main.py}} and continue with the installation steps below.

\Needspace{10\baselineskip}
\section*{Requirements}
\pdfbookmark[1]{Requirements}{chapter-2-s2}
\begingroup
\fontsize{8.7}{11.8}\selectfont
\setlength{\tabcolsep}{5pt}
\renewcommand{\arraystretch}{1.23}
\begin{longtable}{@{}>{\raggedright\arraybackslash}p{\dimexpr 0.34000\linewidth-3.400pt\relax}>{\raggedright\arraybackslash}p{\dimexpr 0.66000\linewidth-6.600pt\relax}@{}}
\rowcolor{navy} \color{white}\bfseries Requirement & \color{white}\bfseries Details \\
\endfirsthead
\rowcolor{navy} \color{white}\bfseries Requirement & \color{white}\bfseries Details \\
\endhead
\endfoot
\rowcolor{pale}
Python & Project baseline: 3.10+. Review environment: Windows with Python 3.14.6. \\
Display & Qt-compatible desktop; minimum window 1200 × 700, initial size 1600 × 950 \\
\rowcolor{pale}
Packages & Install all entries in \code{\seqsplit{requirements.txt}} \\
Source & Keep \code{\seqsplit{core/}}, \code{\seqsplit{gui/}}, \code{\seqsplit{protocols/}}, stylesheet, and demo beside \code{\seqsplit{main.py}} \\
\rowcolor{pale}
Network & Local UDP/TCP sockets; isolated laboratory for external receivers \\
Optional tools & Wireshark/capture reader and an independently configured PDC \\
\end{longtable}
\endgroup

Windows supports the Qt interface and main UDP workflow. GOOSE uses Linux \code{\seqsplit{AF\_PACKET}}; native Windows does not support that publisher. Linux/WSL GOOSE also depends on raw socket permissions and the interface. The core experiment does not require elevating the desktop application.

\Needspace{10\baselineskip}
\section*{Windows installation}
\pdfbookmark[1]{Windows installation}{chapter-2-s3}
\par\noindent\begin{minipage}{\linewidth}
Open PowerShell in the directory containing \code{\seqsplit{main.py}}:
\begin{lstlisting}[title={POWERSHELL}]
py -3 --version
py -3 -m venv .venv
.\.venv\Scripts\python.exe -m pip install -r requirements.txt
.\.venv\Scripts\python.exe -c "import PyQt6, numpy, matplotlib; print('Core imports OK')"
.\.venv\Scripts\python.exe main.py --demo
\end{lstlisting}
\end{minipage}\par
Calling the interpreter directly avoids a PowerShell activation-policy change. Use an existing working environment if you already have one. Do not reuse an environment copied from another OS or machine.

\Needspace{10\baselineskip}
\section*{Linux or WSL installation}
\pdfbookmark[1]{Linux or WSL installation}{chapter-2-s4}
\par\noindent\begin{minipage}{\linewidth}
\begin{lstlisting}[title={BASH}]
python3 --version
python3 -m venv .venv
.venv/bin/python -m pip install -r requirements.txt
.venv/bin/python -c "import PyQt6, numpy, matplotlib; print('Core imports OK')"
.venv/bin/python main.py --demo
\end{lstlisting}
\end{minipage}\par
If environment creation fails, install your distribution's Python virtual-environment package. Qt may also need XCB, OpenGL, EGL, font, and D-Bus libraries.

The bundled \code{\seqsplit{install.sh}} targets Ubuntu 22.04 and creates \code{\seqsplit{venv/}}. It invokes \code{\seqsplit{sudo}\ \seqsplit{apt-get}}, selects Python, installs system packages, and installs requirements. Review it for other distributions. Its system-package block suppresses failures with \code{\seqsplit{||}\ \seqsplit{true}}; a success banner alone does not verify the installation.

\par\noindent\begin{minipage}{\linewidth}
\begin{lstlisting}[title={BASH}]
bash install.sh
venv/bin/python main.py --demo
\end{lstlisting}
\end{minipage}\par
\Needspace{10\baselineskip}
\section*{Display setup}
\pdfbookmark[1]{Display setup}{chapter-2-s5}
On Windows, use a normal desktop session. Under WSL, use an available WSLg display or a configured display server. Preserve working \code{\seqsplit{DISPLAY}} and \code{\seqsplit{WAYLAND\_DISPLAY}} values.

On Linux, \code{\seqsplit{main.py}} sets \code{\seqsplit{QT\_QPA\_PLATFORM=offscreen}} when neither display variable exists, unless a platform was already set. Offscreen mode supports checks but produces no visible window. It is not a headless simulation service.

\Needspace{10\baselineskip}
\section*{First experiment}
\pdfbookmark[1]{First experiment}{chapter-2-s6}
\begin{enumerate}
\item Launch with \code{\seqsplit{--demo}}; it loads a topology without starting simulation.
\item Choose \textbf{Scenarios  >  Clean Baseline}.
\item Click \textbf{Validate}; the demo has a PMU, Local PDC, and attached Threat Agent.
\item Click \textbf{Run} (\code{\seqsplit{F5}}). Check packet totals and proxy status.
\item Inspect magnitude, frequency, and angle. Clean values overlap, apart from possible startup/configuration-record artifacts.
\item Click \textbf{Stop} (\code{\seqsplit{F6}}).
\item Choose \textbf{Scenarios  >  Noisy Sensor}, then run again. Modified magnitude varies around the original.
\item Stop, save the topology, and export incident JSON/CSV from the Integrations export tab.
\end{enumerate}

The demo forwards UDP to \code{\seqsplit{127.0.0.1:4713}}. A Local PDC node represents a destination; it does not launch a PDC receiver. The GUI can display proxy results with no listener. See \hyperref[chapter-9]{Integrations} for receipt verification.

\Needspace{10\baselineskip}
\section*{Installation verification}
\pdfbookmark[1]{Installation verification}{chapter-2-s7}
\par\noindent\begin{minipage}{\linewidth}
Windows:
\begin{lstlisting}[title={POWERSHELL}]
.\.venv\Scripts\python.exe test_regression.py
.\.venv\Scripts\python.exe docs/examples/inspect_frame.py
\end{lstlisting}
\end{minipage}\par
\par\noindent\begin{minipage}{\linewidth}
Linux/WSL:
\begin{lstlisting}[title={BASH}]
.venv/bin/python test_regression.py
.venv/bin/python docs/examples/inspect_frame.py
\end{lstlisting}
\end{minipage}\par
\par\noindent\begin{minipage}{\linewidth}\setlength{\parskip}{5pt}
The regression suite uses offscreen Qt. The documentation example verifies a deterministic transformation without network traffic. See \hyperref[chapter-15]{Testing} for results and coverage limits, then \hyperref[chapter-3]{User guide} to build an experiment.
{\footnotesize\color{muted}\raggedright Source: \href{https://github.com/neerajjayesh/GridSec-project/blob/main/main.py}{entry point} (\code{\seqsplit{main.py}}), \href{https://github.com/neerajjayesh/GridSec-project/blob/main/requirements.txt}{requirements} (\code{\seqsplit{requirements.txt}}), \href{https://github.com/neerajjayesh/GridSec-project/blob/main/install.sh}{installer} (\code{\seqsplit{install.sh}}), \href{https://github.com/neerajjayesh/GridSec-project/blob/main/mitm-demo-topology.json}{demo} (\code{\seqsplit{mitm-demo-topology.json}}).\par}
\end{minipage}\par

